\section{The exact interval criterion}\label{sec:criterion}

\subsection{Notation and the elementary endpoint lemma}
Let
\[
\begin{aligned}
f(n)&=\lfloor n/2\rfloor+\lfloor n/3\rfloor-\lfloor n/6\rfloor,
& m(n)&=\lfloor n/6\rfloor,\\
H(n)&=\lceil n/2\rceil,
& M(n)&=\lfloor n/3\rfloor-\lfloor n/6\rfloor+1.
\end{aligned}
\]
Let $O_U$ be the odd integers in $[U]$, in any fixed total order (one may use decreasing $\phi(v)/v$). Write $P_k$ for its first $k$ members.

The retained odd count has the useful identity
\[
 M(n)=\left\lfloor\frac{n+3}{6}\right\rfloor+1.
\]
Any set of $M(n)$ distinct odd integers in $[n]$, for $n\ge6$, contains a coprime pair. Here is a direct proof. If the set contains $1$, choose any other member; $M(n)\ge2$. Otherwise map every member $a$ to $\lfloor(a+3)/6\rfloor$. The image lies in $\{1,\ldots,\lfloor(n+3)/6\rfloor\}$, so two distinct members lie in one fibre. Each fibre consists of three consecutive odd integers, of the form $\{6k+3,6k+5,6k+7\}$. Their pairwise differences are $2$ or $4$; a common divisor of two odd members must therefore be $1$.

This is the entire endpoint input in the formal proof: \texttt{oddEndpointPrinciple} in \path{Erdos883AKEndpoints.lean}. The historical filename is retained, but no Ahlswede--Khachatrian theorem is assumed. The argument replaces the external endpoint input used in the earlier informal manuscript and in the cited asymptotic architecture.

Choose any ordered finite list of divisor coordinates $c_1,\ldots,c_t$. The $s$-signature of $v$ is $(1_{c_i\mid v})_{1\le i\le s}$. The graph-theoretic criterion does not require these coordinates to be prime. Set $h_0=0$ and $h_s=2^s+1$ for $s\ge1$.
For integers $6\le L\le U$, define:
\[
\begin{aligned}
e_s(k;L,U)
&=\min_{\substack{u,v\in P_k,\ u\ne v\\\text{equal $s$-signature}}}
\#\{w\in[L]:w\text{ even},\ \gcd(w,u)=\gcd(w,v)=1\};\\
r_s(k;L,U)
&=\text{the same minimum with all positive }w\le L.
\end{aligned}
\]
If there is no such pair, set the corresponding minimum to $+\infty$. In Lean, \texttt{ExactIntervalCertificate} expresses the corresponding lower bound pointwise for every eligible pair, so it needs no infinite value or numerical sentinel.

\subsection*{Exact interval criterion}
Suppose for every integer
\[
0\le b\le H(U)-M(U),\qquad 1\le j\le m(U),
\]
there exist an $s\in\{0,\ldots,t\}$ and $k\in\{0,\ldots,H(U)\}$ such that
\begin{equation}
2\max(0,H(U)-k-b)+h_s<j, \tag{1}
\end{equation}
and at least one of
\begin{align}
e_s(k;L,U)&\ge b+j, \tag{2E}\\
r_s(k;L,U)&\ge U-f(U)+m(U)+j. \tag{2R}
\end{align}
Then for every integer $n\in[L,U]$ and every $A\subseteq[n]$ with $|A|>f(n)$, $G(A)$ contains every odd cycle $C_{2\ell+1}$, $1\le\ell\le\lfloor n/6\rfloor$.

\subsection*{Proof}
Fix such $n,A$ and put $b=\#(\text{even integers in }[n]\text{ outside }A)$. There are at least $M(n)+b$ selected odds. Retain the first $M(n)$ of these in the fixed $O_U$ order, calling the resulting set $Y_0$. For every $k$,
\begin{equation}
|Y_0\setminus P_k|\le\max(0,H(U)-k-b). \tag{3}
\end{equation}
Indeed, if a retained odd lies after position $k$, every discarded selected odd lies after $k$ as well; there are at least $b$ discarded odds. Counting inside $O_U$ proves~(3). If no retained odd lies after $k$ the assertion is immediate.

By the elementary endpoint lemma, choose coprime distinct $a,a'$ in $Y_0$. Since $M(n)\ge m(n)+1$, for each $1\le\ell\le m(n)$ choose $\ell-1$ other members of $Y_0$. Sort these middle vertices lexicographically by their full $t$-signatures, then place $a$ first and $a'$ last. Denote the result $o_0=a,\ldots,o_\ell=a'$. This is the construction in \texttt{exists\_divisor\_skeleton}.

This one order has, simultaneously for every $1\le s\le t$, at most $h_s$ gaps joining different $s$-signatures: within the sorted middle, every prefix-signature occupies one block and there are at most $2^s-1$ transitions; adjoining the two prescribed endpoints adds at most two. This argument also covers an empty middle. For $s=0$ there are no transitions. At most twice the number in~(3) of the gaps touch an odd outside $P_k$.

Let $R=A$ minus the skeleton and let
\[
D_i=\{w\in R:\gcd(w,o_{i-1})=\gcd(w,o_i)=1\}.
\]
For each rank $j\le\ell$ take $s,k$ from the criterion. Every gap except at most $2\max(0,H(U)-k-b)+h_s$ has both endpoints in $P_k$ and equal $s$-signature.

If (2E) holds, that gap has at least $b+j$ common even neighbors in $[L]$. At most $b$ are missing from $A$, and no even vertex is in the odd skeleton, so $|D_i|\ge j$.

If (2R) holds, it has at least $U-f(U)+m(U)+j$ common neighbors in $[L]$. Removing vertices outside $A$ and in the skeleton costs at most
\[
(n-|A|)+(\ell+1)\le n-f(n)+\ell\le U-f(U)+m(U),
\]
so again $|D_i|\ge j$. The last inequality follows from monotonicity of $n-f(n)+m(n)$; writing $n=6q+r$, it equals $3q+(0,1,1,1,1,2)_r$.
Also $b\le H(n)-M(n)\le H(U)-M(U)$, since $H(n)-M(n)=2q+(-1,0,0,0,0,1)_r$ is nondecreasing.

Consequently fewer than $j$ lists $D_i$ have size less than $j$. Sort the $\ell$ lists by increasing cardinality. The $j$th list has size at least $j$, so greedy selection supplies distinct representatives $w_i$. They all lie outside the skeleton. Hence
\[
o_0,w_1,o_1,\ldots,w_\ell,o_\ell,o_0
\]
is a simple coprime cycle of length $2\ell+1$, closing via $\gcd(a,a')=1$. Different ranks may use different $s,k$ or (2E)/(2R), because each bounds the same actual resource lists $D_i$. This proves the criterion.

\section{Arithmetic resources}\label{sec:resources}

\subsection*{Notation}
$\rho(v)=\phi(v)/v$. Throughout this part, endpoint pairs $u,v$ are odd integers.
For the initial tail estimates, specialize the coordinates to the first $s$ odd primes $p_1=3,\ldots,p_s$, and write $p=p_{s+1}$. The uniform construction below will instead use a specified coherent family of divisor-coordinate lists.
For $s=0$, set $p_0=2$. The $0$-signature is empty, $p=p_1=3$, and the tail primes $r>p_0$ are exactly the odd primes; all endpoint pairs remain odd.
Let $U$ be an integer upper bound on every skeleton vertex and let $L\le U$.
Let
\[
\begin{aligned}
H&=\lceil U/2\rceil,& q_L&=\lfloor L/2\rfloor,& m&=\lfloor U/6\rfloor,\\
f&=\lfloor U/2\rfloor+\lfloor U/3\rfloor-\lfloor U/6\rfloor,
& M&=f-\lfloor U/2\rfloor+1.
\end{aligned}
\]
For $0<z\le1$ define the strict finite profile count
\[
R_U(z)=\#\{\text{odd }v\le U:\rho(v)<z\}.
\]
Strict inequality ensures that every remaining vertex has profile $\ge z$.
Write $h_0=0$ and $h_s=2^s+1$ for $s\ge1$.

\subsection*{1. Exact adaptive tail factor}
Let $a$ be the greatest nonnegative integer with $p^a\le U$ and put
\[
\eta_s(U)=\left(\frac{p-1}{p}\right)^a.
\]
Every $v\le U$ has at most $a$ distinct prime divisors exceeding $p_s$. Each corresponding totient factor is at least $(p-1)/p$. Therefore
\[
\prod_{\substack{r\mid v,\ r>p_s\\r\text{ prime}}}\left(1-\frac1r\right)\ge\eta_s(U).
\]
All ingredients are integers or rational numbers; logarithms are unnecessary for computing $a$.

If odd integers $u,v\le U$ share their first $s$ prime-divisibility signatures, every prime factor of $v$ absent from $u$ exceeds $p_s$. Hence
\[
\begin{aligned}
\rho(\operatorname{lcm}(u,v))
&=\rho(u)\prod_{\substack{r\mid v,\ r\nmid u\\r\text{ prime}}}\left(1-\frac1r\right)\\
&\ge\rho(u)\eta_s(U).
\end{aligned}
\]
Interchanging $u,v$ proves the slightly stronger bound
\[
\rho(\operatorname{lcm}(u,v))\ge\max(\rho(u),\rho(v))\eta_s(U).
\]
There is also the signature-free inequality
\[
 \rho(\operatorname{lcm}(u,v))\ge\rho(u)\rho(v),
\]
because the Euler factors at common prime divisors appear twice on the right and once on the left. Thus when both profiles are at least $z\ge0$ the useful combined bound is
\[
 B_s(z):=\max\{z^2,z\eta_s(U)\}\le\rho(\operatorname{lcm}(u,v)).
\]
The final finite adaptive certificates use this combined bound.

One may sharpen $\eta_s(U)$ without changing the proof. Let $r$ be maximal such that the product of the first $r$ primes after $p_s$ is $\le U$. No integer $\le U$ has more than $r$ distinct tail prime factors. Among any $r$ such primes the smallest $r$ give the smallest totient-factor product; fewer factors give no smaller a product. Thus their factor product is a valid sharper rational $\eta_s(U)$. The formal \texttt{SharpTailCertificate} packages this argument as a finite set $Q$ and threshold $q$: all odd primes below $q$ belong either to the signature list or to $Q$, every $x\in Q$ satisfies $2\le x\le q$, and
\[
 U<q\prod_{x\in Q}x,\qquad
 \eta_s(U)=\prod_{x\in Q}(1-1/x).
\]
A product-barrier lemma and an Euler-factor comparison prove the resulting tail bound. In the supplied certificates $Q$ consists of the initial primes beyond the signature prefix.

\subsection*{2. Exact common-neighbor discrepancy}
For odd squarefree $D$ with $w=\omega(D)\ge1$ and integer $X\ge0$,
\[
\begin{aligned}
C_D(X)&=\#\{1\le a\le X:\gcd(a,D)=1\}\\
&=\sum_{d\mid D}\mu(d)\lfloor X/d\rfloor.
\end{aligned}
\]
Comparing with $X\phi(D)/D$, only the positive-$\mu$ terms can contribute negatively. There are $2^{w-1}$ such divisors, and $d=1$ has zero fractional part. Thus
\[
C_D(X)\ge X\rho(D)-(2^{w-1}-1).
\]
For $D=1$ the equality $C_1(X)=X$ holds.

Define $W(U)$ as the greatest $w$ such that the product of the first $w$ odd primes is $\le U^2$, and $\Delta(U)=2^{W(U)-1}-1$ (take zero if $W(U)=0$).
For $D=\operatorname{rad}(uv)$, $D\le uv\le U^2$, so the previous inequality gives
\[
\begin{aligned}
\text{common even neighbors in }[L]&\ge q_L\rho(D)-\Delta(U),\\
\text{common whole-pool neighbors in }[L]&\ge L\rho(D)-\Delta(U).
\end{aligned}
\]
This does not require $D$ to divide $L$ and remains valid if an endpoint is $1$.

\subsection*{3. Exact finite criterion using only a univariate profile}
For every required pair $0\le b\le H-M$ and $1\le j\le m$, suppose one can choose a signature prefix $s$ and rational $z$ with
\begin{equation}
2\max(0,R_U(z)-b)+h_s<j, \tag{A}
\end{equation}
and either
\begin{equation}
q_LB_s(z)-\Delta(U)\ge b+j, \tag{$B_E$}
\end{equation}
or
\begin{equation}
LB_s(z)-\Delta(U)\ge U-f+m+j. \tag{$B_R$}
\end{equation}
Then the all-subsets conclusion of the finite interval criterion holds for every $n\in[L,U]$.

\noindent\textit{Proof:} retain the best $M(n)$ selected odds as before. Among the $R_U(z)$ low profile odds, at most $\max(0,R_U(z)-b)$ can remain in that retained set, because at least $b$ selected odds worse than any retained low-profile odd were discarded. Choose a single full coordinate list containing all prefixes used in the certificate, sort the middle of the skeleton by its full signature, and adjoin its coprime endpoints. For any selected prefix $s$, all but the number in (A) of its gaps have equal signatures and both endpoint profiles $\ge z$. The tail and discrepancy bounds give at least $j$ resources for each such gap. The ordered-Hall proof applies unchanged.

Different ranks may select different prefix lengths. A single ordering by the full list has all shorter prefix signatures contiguous in the middle, so there is no consistency problem.

For the finite implementation, define integer lower degrees
\[
 E_s(v)=\lfloor q_LB_s(\rho(v))-\Delta(U)\rfloor_+,
 \qquad
 Q_s(v)=\lfloor LB_s(\rho(v))-\Delta(U)\rfloor_+,
\]
where $\lfloor y\rfloor_+=\max(0,\lfloor y\rfloor)$. Both are nondecreasing functions of the same profile $\rho(v)$. For each $(b,j)$ it suffices to select $s$ with $h_s<j$ and verify either
\[
 \#\{v\text{ odd}\le U:E_s(v)<b+j\}
 \le b+\left\lfloor\frac{j-h_s-1}{2}\right\rfloor,
\]
or
\[
 \#\{v\text{ odd}\le U:Q_s(v)<U-f+m+j\}
 \le b+\left\lfloor\frac{j-h_s-1}{2}\right\rfloor.
\]
One common decreasing-profile order turns each good-degree set into an initial segment. Its complementary low count gives the strict exceptional-gap budget, while the degree threshold gives the corresponding common-neighbor bound. This is \texttt{DegreeIntervalCertificate}; its reduction to the exact criterion is proved by \texttt{exact\_interval\_of\_degree\_certificate} in \path{Erdos883DegreeCriterion.lean}.

\subsection*{4. A fully uniform tail bound with $O(\sqrt U)$ transitions}
For every integer $U>1024$ choose the least $s$ with $U\le4^s$. Thus $s\ge6$ and $4^{s-1}<U$, whence $2^s<2\sqrt U$.

Here the precise coordinate family matters. Let $\mathcal P$ be the checked list of the first $511$ odd primes, ending at $3671$. Define $\mathcal C_s$ to be the first $s$ entries of $\mathcal P$ when $s\le511$; for $s>511$, append $s-511$ consecutive odd integers beginning with $3673$. These appended coordinates need not be prime. The family has length $s$ and is coherent under taking prefixes. For $s\ge511$, it contains every odd prime below $2s+2651$. The source checks the coverage of the finite list by certifying that omitted integers in $[3,3671]$ have proper divisors.

If $u,v$ share their $\mathcal C_s$-signatures, every prime factor of $v$ absent from $u$ is at least a suitable threshold $p$. Let $a$ count these missing prime factors; then $p^a\le v\le U\le4^s$.

For $s\ge511$, take $p=2s+2651$, so $p\ge2s+3\ge1024$. If $10a\ge p$, then
\[
 4^s=2^{2s}<2^{10a}\le p^a,
\]
a contradiction. Therefore $10a<p$, and Bernoulli's inequality gives
\[
 \prod_{\substack{r\mid v,\ r\nmid u\\r\text{ prime}}}\left(1-\frac1r\right)
 \ge(1-1/p)^a\ge1-a/p>\frac9{10}.
\]
No estimate for the growth of primes is needed.

For $9\le s<511$, let $p$ be entry $s+1$ of $\mathcal P$. The $502$ kernel-checked integer inequalities
\[
 4^s<p^{\lceil p/10\rceil}
\]
again imply $10a<p$. For $s=6,7,8$, separate exact product barriers give the same strict Euler-product bound: at scale $6$ there are at most two missing primes, and if $19$ occurs the other is at least $23$; at scale $7$ there are at most two, all at least $23$; at scale $8$ there are at most three, all at least $29$. The inequalities $(18/19)(22/23)>9/10$, $(22/23)^2>9/10$, and $(28/29)^3>9/10$ complete these cases.

Consequently,
\[
 h_s=2^s+1<2\sqrt U+1,\qquad
 \rho(\operatorname{lcm}(u,v))>\frac9{10}\max(\rho(u),\rho(v))
\]
whenever $u,v\le U$ have equal full signatures. The declarations are \texttt{adaptiveTailCoordinates}, \texttt{uniformTailScale}, and \texttt{totientDensity\_lcm\_gt\_nine\_tenths\_uniform} in \path{Erdos883FiniteTailThresholds.lean}, with the elementary large-scale argument in \path{Erdos883UniformTail.lean} and the three small cases in \path{Erdos883SmallTailCertificates.lean}.

A useful uniform discrepancy bound is also elementary. For odd squarefree $D$,
\[
\frac{(2^{\omega(D)})^4}{D}
=\prod_{p\mid D}\frac{16}{p}
\le\frac{16^5}{3\cdot5\cdot7\cdot11\cdot13}<81,
\]
since factors from primes $\ge17$ are $\le1$. Hence $2^{\omega(D)}<3D^{1/4}$.
For $D\le U^2$, the actual discrepancy from the discrepancy bound is strictly less than $(3/2)\sqrt U$. Thus $(B_E)/(B_R)$ may use loss $(3/2)\sqrt U$ and factor $9/10$, without computing $W(U)$.

This proves a universal numerical reduction for all $U>1024$:
\[
2\max(0,R_U(z)-b)+2^s+1<j,
\]
and either
\[
\frac9{10}q_Lz-\frac32\sqrt U\ge b+j,
\]
or
\[
\frac9{10}Lz-\frac32\sqrt U\ge U-f+m+j
\]
are sufficient. The exact integer $2^s+1$ is preferable to its displayed $2\sqrt U+1$ upper bound. Irrational square roots may be rounded upward in the subtracted error if a rational-only certificate is desired.

